\documentclass[10pt,a4paper]{article}

% ---------- Codificación e idioma ----------
\usepackage[utf8]{inputenc}
\usepackage[T1]{fontenc}            % acentos y guiones bien en el PDF
\usepackage[english]{babel}
\DeclareUnicodeCharacter{2212}{-}   % el signo menos Unicode no rompe la compilación

% ---------- Página ----------
\usepackage[a4paper,margin=2cm]{geometry}   % 2 cm: límite de 2 páginas
\usepackage{microtype}              % mejor reparto de espacios, ahorra líneas

% ---------- Matemáticas ----------
\usepackage{amsmath,amssymb,amsthm}

% ---------- Unidades y química ----------
\usepackage{siunitx}
\sisetup{per-mode=symbol}           % kcal/mol en vez de kcal mol^-1
\DeclareSIUnit\calorie{cal}         % ¡usar \kilo\calorie, no \kilo\cal!
\DeclareSIUnit\atm{atm}
\DeclareSIUnit\hartree{\ensuremath{E_\mathrm{h}}}
\usepackage[version=4]{mhchem}

% ---------- Figuras y tablas ----------
\usepackage{graphicx}
\usepackage{booktabs}
\usepackage{float}
\usepackage[font=small,labelfont=bf]{caption}
\usepackage{tikz}
\usepackage{pgfplots}
\pgfplotsset{compat=1.18}

% ---------- Otros ----------
\usepackage[normalem]{ulem}         % normalem: \emph sigue siendo cursiva

% ---------- Compactación (para el límite de 2 páginas) ----------
\usepackage{titling}                % título más compacto
\setlength{\droptitle}{-1.6cm}
\pretitle{\begin{center}\large\bfseries}\posttitle{\par\end{center}\vspace{-0.6em}}
\preauthor{\begin{center}\normalsize}\postauthor{\par\end{center}\vspace{-0.9em}}
\predate{\begin{center}\small}\postdate{\par\end{center}\vspace{-1.2em}}
\usepackage{titlesec}               % secciones con menos espacio alrededor
\titleformat{\section}{\large\bfseries}{\thesection}{0.6em}{}
\titlespacing*{\section}{0pt}{1.1ex plus 0.3ex}{0.5ex}
\usepackage[shortlabels]{enumitem}  % listas sin espacios extra
\setlist{nosep, leftmargin=*}
\setlength{\intextsep}{6pt plus 2pt minus 2pt}       % espacio alrededor de figuras/tablas [H]
\captionsetup{skip=4pt}
\setlength{\parskip}{0pt}
\usepackage{etoolbox}               % bibliografía más pequeña y sin huecos
\apptocmd{\thebibliography}{\footnotesize\setlength{\itemsep}{0pt}\setlength{\parsep}{0pt}}{}{}

% ---------- Referencias (siempre al final y en este orden) ----------
\usepackage[hidelinks]{hyperref}
\usepackage[capitalise,noabbrev]{cleveref}



\title{Conformational Analysis of 1,2-Dichloroethane:\\
       the Trans--Gauche Equilibrium at the HF and B3LYP Levels}
\author{Martín Llorens Giménez\thanks{\textbf{Declaration of AI use:} An AI assistant (Claude Opus 5.5, Anthropic) was used for guidance to help structure the text, which the author then wrote independently. AI was mainly used to revise English expressions and grammar, and to write the LaTeX code for figures and tables. All calculations, data extraction, and final decisions on the content are the author's own.}}
\date{October 2026}


\begin{document}

\maketitle

%Attention:
% - 2 pages maximum
% - Numbering of figures and tables, inclusion of figure/table captions
% - Units!
% - Use correct number of significant figures (we will discuss this)


\section{Introduction} %What have i done?

Rotation about a C--C single bond is often described as nearly free. However,
substituents turn the torsional potential into a landscape of minima and barriers
that determines which conformers are populated at a given temperature.
1,2-Dichloroethane (\ce{ClCH2CH2Cl}, DCE) is the textbook example: it exists as a
trans conformer (Cl--C--C--Cl dihedral angle of \ang{180}) and a gauche conformer
(dihedral angle close to \ang{70}). Its relative stability results from a balance
between the repulsion of the two C--Cl bonds, which favours the trans form, and
stabilising hyperconjugative interactions (the \emph{gauche effect}), which favour
the gauche form~\cite{wiberg1987}. A correct description of this balance requires
electron correlation, which makes DCE a demanding test for the Hartree--Fock (HF)
method, which neglects it.

In this work, which corresponds to Exercise~6 of the practical session, the trans
and gauche minima were optimized at the HF/6-31G* level, starting from the
approximate geometries provided by a relaxed scan of the Cl--C--C--Cl dihedral angle
(Exercise~5). Harmonic frequency calculations confirmed both structures as true
minima and provided the thermochemical data needed to obtain $\Delta E$, $\Delta H$,
$\Delta G$ and the trans~$\rightleftharpoons$~gauche equilibrium constant at
\SI{298}{\kelvin}. Going beyond the scope of the exercise, the optimizations and
frequency calculations were repeated with the B3LYP functional, to test whether the
deviation from experiment arises from the lack of electron correlation in HF.

\section{Methods} %How did i do it?

All calculations were performed with Gaussian~09~\cite{gaussian09}, using the
Hartree--Fock (HF) method and the B3LYP functional with the 6-31G* basis set.

First, the potential energy curve for the internal rotation around the C--C bond
was obtained through a relaxed scan of the Cl--C--C--Cl dihedral angle at the
HF/6-31G* level. The dihedral angle was varied from \ang{180} to \ang{0} in steps
of \ang{20}, and all other internal coordinates were optimized at each point.

Then, the geometries of the trans and gauche conformers were optimized in internal
coordinates (Z-matrix), starting from dihedral angles of \ang{180} and \ang{70},
respectively, close to the minima found in the scan. Harmonic vibrational
frequencies were computed for each optimized geometry with two aims: to confirm
that the structures are true minima (no imaginary frequencies) and to obtain the
thermodynamic quantities ($E$, $H$ and $G$) at \SI{298.15}{\kelvin} and 1~atm,
within the ideal gas, rigid rotor and harmonic oscillator approximations. All
optimizations and frequency calculations were repeated with the B3LYP functional.

Finally, the trans~$\rightleftharpoons$~gauche equilibrium constant was calculated
as $K = 2\exp(-\Delta G/RT)$, where the factor of 2 accounts for the two
enantiomeric gauche conformers, which are energetically equivalent.

\section{Results and discussion} %What are the results?

The potential energy curve (\cref{fig:scan}) presents two minima: the trans conformer
(\ang{180}), which is the global minimum, and the gauche conformer, which the full
optimization places at \ang{70.5} (\cref{tab:abs}), 1.91~kcal/mol above. The maxima correspond to
eclipsed conformations: at \ang{120} (5.05~kcal/mol) each Cl eclipses an H, and at
\ang{0} (9.98~kcal/mol) the two Cl atoms are eclipsed. The optimized gauche minimum
lies below the two adjacent scan points, because a \ang{20} step is too coarse to
resolve the bottom of the well. At both levels of theory, all frequencies are real
(the lowest, between 114 and \SI[per-mode=power]{128}{\per\centi\metre}, corresponds to the torsion
about the C--C bond), which confirms that both structures are true minima.

\begin{figure}[H]
  \centering
  % Proyección de Newman como "pic": #1 = ángulo del Cl de detrás (el de delante está a 90°)
  \tikzset{
    pics/newman/.style={code={
      \fill[white] (0,0) circle (0.62);                                     % fondo: tapa la rejilla
      \draw[thick] (0,0) circle (0.25);                                     % carbono de detrás
      \foreach \a in {0,120,240} \draw[thick] (#1+\a:0.25) -- (#1+\a:0.48); % enlaces de detrás
      \foreach \a in {90,210,330} \draw[thick] (0,0) -- (\a:0.48);          % enlaces de delante
      \node[font=\scriptsize, text=green!40!black, inner sep=0pt] at (#1:0.64) {Cl};
      \node[font=\scriptsize, text=green!40!black, inner sep=0pt] at (90:0.64) {Cl};
    }}
  }
  \begin{tikzpicture}
    \begin{axis}[
        width=\textwidth, height=4.6cm,
        xlabel={Cl--C--C--Cl dihedral angle (\si{\degree})},
        ylabel={$\Delta E$ (kcal/mol)},
        xmin=-5, xmax=193, ymin=-0.8, ymax=11,
        xtick={0,20,...,180}, ytick={0,2,...,10},
        grid=major, grid style={gray!25},
        tick label style={font=\footnotesize}, label style={font=\small},
        legend style={font=\footnotesize, at={(0.015,0.03)}, anchor=south west,
                      draw=none, fill=white, fill opacity=0.85, text opacity=1},
        clip=false,
        % Marca extra en el eje X para el diedro gauche optimizado (sin línea de rejilla)
        extra x ticks={70.5},
        extra x tick labels={\color{red!60!black}70.5},
        extra x tick style={grid=none, tick label style={font=\scriptsize}},
      ]
      \addplot[thick, blue!70!black, mark=*, mark size=1.6pt]
        table[x=dihedral, y=dE_kcal] {scan_hf.dat};
      \addlegendentry{Relaxed scan}
      \addplot[only marks, mark=diamond*, mark size=3.2pt, red!75!black]
        table[x=dihedral, y=dE_kcal] {minima_hf.dat};
      \addlegendentry{Optimized minima}
      % Puntos donde irán las proyecciones (en coordenadas de la gráfica)
      \coordinate (pg) at (axis cs:70.5,6.6);
      \coordinate (pt) at (axis cs:179.9999,5.0);
      % Conectores: desde el eje X, pasando por el mínimo, hasta su proyección (no más arriba)
      \draw[densely dashed, red!60!black, thin] (axis cs:70.5,\pgfkeysvalueof{/pgfplots/ymin}) -- (pg);
      \draw[densely dashed, red!60!black, thin] (axis cs:179.9999,\pgfkeysvalueof{/pgfplots/ymin}) -- (pt);
    \end{axis}
    % Proyecciones y etiquetas, fuera del axis: así se dibujan en cm reales, sin deformarse
    \pic at (pg) {newman=160.5};
    \pic at (pt) {newman=270};
    \node[font=\footnotesize, anchor=west] at ([xshift=0.85cm]pg) {gauche (\ang{70.5})};
    \node[font=\footnotesize, anchor=east] at ([xshift=-0.85cm]pt) {trans (\ang{180})};
  \end{tikzpicture}
  \caption{Potential energy curve for the internal rotation of 1,2-dichloroethane at the
    HF/6-31G* level (relaxed scan, energies relative to the trans conformer). Diamonds mark
    the fully optimized minima, shown above as Newman projections along the C--C bond.}
  \label{fig:scan}
\end{figure}

\Cref{tab:abs} lists the geometries and energies of the two conformers. Both methods give
essentially the same geometry: the gauche dihedral angle barely changes between HF
(\ang{70.5}) and B3LYP (\ang{69.6}), and both values agree with the electron diffraction
value of $72.4 \pm 3.1^{\circ}$~\cite{wiberg1987}. The absolute HF and B3LYP energies differ by
about 2~$E_\mathrm{h}$, but this difference has no physical meaning, because each method has
its own energy reference; only differences computed with the same method are comparable,
which are those in \cref{tab:delta}. The HF/6-31G* energies coincide with those reported by
Wiberg and Murcko~\cite{wiberg1987}, which validates the procedure. Finally, the trans conformer has no dipole moment, as
required by its centre of inversion, whereas the gauche conformer has a sizeable dipole
moment ($\approx$3~D), slightly smaller with B3LYP, consistent with the well-known tendency
of HF to overestimate dipole moments.

\begin{table}[H]
  \centering
  \caption{Optimized dihedral angles, electronic energies ($E$), enthalpies ($H$) and Gibbs
    free energies ($G$) at \SI{298.15}{\kelvin}, and dipole moments ($\mu$) of the trans and
    gauche conformers.}
  \label{tab:abs}
  \footnotesize\setlength{\tabcolsep}{4.5pt}
  \begin{tabular}{ll
                  S[table-format=3.1]
                  S[table-format=-3.6]
                  S[table-format=-3.6]
                  S[table-format=-3.6]
                  S[table-format=1.2]}
    \toprule
    Level & Conformer & {Dihedral (\si{\degree})} & {$E$ ($E_\mathrm{h}$)}
          & {$H$ ($E_\mathrm{h}$)} & {$G$ ($E_\mathrm{h}$)} & {$\mu$ (D)} \\
    \midrule
    HF/6-31G*    & trans  & 180.0 & -997.030936 & -996.962335 & -996.996809 & 0.00 \\
                 & gauche &  70.5 & -997.027890 & -996.959483 & -996.993863 & 3.14 \\
    \addlinespace
    B3LYP/6-31G* & trans  & 180.0 & -999.018999 & -998.954516 & -998.989373 & 0.00 \\
                 & gauche &  69.6 & -999.016295 & -998.952027 & -998.986742 & 2.93 \\
    \bottomrule
  \end{tabular}
\end{table}

As shown in \cref{tab:delta}, $\Delta E$, $\Delta H$ and $\Delta G$ are very similar: the
thermal corrections change the difference by less than 0.15~kcal/mol and the entropy
difference is almost zero ($\Delta S \approx -0.2$~cal\,mol$^{-1}$\,K$^{-1}$), so the
equilibrium is governed by the electronic energy. At the HF level, $K = 0.088$, which
corresponds to 8\% of the gauche conformer at \SI{298}{\kelvin}. However, HF overestimates the
enthalpy difference by about 0.7~kcal/mol relative to the experimental value
($1.1 \pm 0.1$~kcal/mol, mostly from the temperature dependence of gas-phase
measurements)~\cite{wiberg1987} and predicts a gauche population of about one
third of the experimental estimate ($\approx$24\%). This discrepancy illustrates the
exponential dependence of $K$ on $\Delta G$: an error of 0.75~kcal/mol in $\Delta G$ reduces
$K$ by a factor of about 3.5. This is why computational chemistry aims at so-called
\emph{chemical accuracy} (about 1~kcal/mol).

\begin{table}[H]
  \centering
  \caption{Gauche $-$ trans energy differences, equilibrium constant and
    gauche population at \SI{298.15}{\kelvin}.}
  \label{tab:delta}
  \small
  \begin{tabular}{l
                  S[table-format=1.2]
                  S[table-format=1.2]
                  S[table-format=1.2]
                  S[table-format=1.3]
                  S[table-format=2.0]}
    \toprule
    Level & {$\Delta E$ (kcal/mol)} & {$\Delta H$ (kcal/mol)} & {$\Delta G$ (kcal/mol)}
          & {$K$} & {Gauche (\%)} \\
    \midrule
    HF/6-31G*    & 1.91 & 1.79 & 1.85 & 0.088 & 8  \\
    B3LYP/6-31G* & 1.70 & 1.56 & 1.65 & 0.12  & 11 \\
    Experiment   & {--} & {$1.1 \pm 0.1$\textsuperscript{a}} & {--}
                 & {$\approx 0.31$\textsuperscript{b}} & {$\approx 24$\textsuperscript{b}} \\
    \bottomrule
  \end{tabular}

  \smallskip
  {\footnotesize \textsuperscript{a}Gas-phase value, mostly from temperature-dependence
   measurements~\cite{wiberg1987}. \textsuperscript{b}Estimated from the experimental value,
   assuming $\Delta G \approx \Delta H$.}
\end{table}

To test whether this deviation is due to the lack of electron correlation, the
calculations were repeated with the B3LYP functional. Including correlation reduces
$\Delta H$ to 1.56~kcal/mol and raises the gauche population to 11\%: the result
improves in the expected direction, but corrects only about a third of the error.
Electron correlation is therefore necessary, but not sufficient at this level of
theory. The main limitation is not the basis set: Wiberg and Murcko~\cite{wiberg1987}
showed that adding diffuse functions barely changes the result (from 1.60 to
1.54~kcal/mol at the MP3 level). The problem lies in the treatment of correlation:
B3LYP is approximate and does not describe dispersion, which would favour the more
compact gauche conformer, and even with explicit correlation (MP3, MP4) the difference
remains around 1.5~kcal/mol.

\section{Conclusion} %What can i conclude?

\begin{enumerate}
  \item 1,2-Dichloroethane has two stable conformers, trans and gauche, both confirmed
    as true minima by frequency calculations. At the HF/6-31G* level, the trans conformer
    is more stable, with $\Delta G = 1.85$~kcal/mol, which gives $K = 0.088$ and a gauche
    population of 8\% at \SI{298}{\kelvin}.
  \item The three energy differences ($\Delta E$, $\Delta H$ and $\Delta G$) are
    practically equal, so the conformational equilibrium is governed by the electronic
    energy rather than by entropic or thermal effects.
  \item HF overestimates the stability of the trans conformer by about 0.7~kcal/mol and,
    given the exponential dependence of $K$ on $\Delta G$, underestimates the gauche
    population by a factor of about 3. Including electron correlation with B3LYP improves
    the result but corrects only about a third of the error: correlation is necessary but
    not sufficient, and a quantitative description requires a more complete treatment of
    electron correlation, including dispersion.
  \item Finally, all results refer to the gas phase. Since the trans conformer has no
    dipole moment ($\mu = 0$ by symmetry) whereas the gauche conformer does
    ($\mu \approx 3$~D), a polar solvent is expected to stabilize the gauche conformer
    and shift the equilibrium. This could be studied with an implicit solvent model.
\end{enumerate}
\bibliographystyle{unsrt}
\bibliography{bibliografia}

\end{document}
